\section{Classical information and operator-Schmidt bounds}

\begin{definition}[Joint probability distribution and marginals]
    \label{def:finite_joint_distribution_marginals}
    \lean{Matrix.IsJointDistribution,
        Matrix.rowMarginal,
        Matrix.columnMarginal}
    \leanok
    Let $X$ and $Y$ be finite sets. A matrix
    $P=(P_{x,y})_{x\in X,y\in Y}$ is a joint probability distribution if
    $P_{x,y}\geq0$ for every $x\in X$ and $y\in Y$, and
    \begin{align}
        \sum_{x\in X}\sum_{y\in Y}P_{x,y}
         & =1.
        \notag
    \end{align}
    Its row and column marginals are respectively
    \begin{align}
        p_x & =\sum_{y\in Y}P_{x,y},
        \notag                       \\
        q_y & =\sum_{x\in X}P_{x,y}.
        \notag
    \end{align}
\end{definition}

\begin{definition}[Entropy and classical mutual information]
    \label{def:finite_classical_mutual_information}
    \lean{Entropy.probabilityEntropy,
        Entropy.classicalMutualInformation}
    \leanok
    \uses{def:finite_joint_distribution_marginals}
    For $t>0$, set $h(t)=-t\log t$, and set $h(0)=0$. The entropy of a
    probability distribution $a=(a_z)_{z\in Z}$ on a finite set is
    $H(a)=\sum_{z\in Z}h(a_z)$. For a joint probability distribution $P$
    with row and column marginals $p$ and $q$, put
    \begin{align}
        H(P)     & =\sum_{x\in X}\sum_{y\in Y}h(P_{x,y}),
        \notag                                            \\
        I(X:Y)_P & =H(p)+H(q)-H(P).
        \notag
    \end{align}
\end{definition}

\begin{lemma}[Entropy is bounded by the logarithm of the support cardinality]
    \label{lem:probability_entropy_le_log_card_support}
    \lean{Entropy.probabilityEntropy_le_log_card_support}
    \leanok
    \uses{def:finite_classical_mutual_information}
    Let $p=(p_z)_{z\in Z}$ be a probability distribution on a finite set $Z$.
    Define $h(t)=-t\log t$ for $t>0$ and $h(0)=0$. Then
    \begin{align}
        \sum_{z\in Z}h(p_z)
         & \leq
        \log\left\lvert\{z\in Z:p_z\neq0\}\right\rvert.
        \notag
    \end{align}
\end{lemma}
\begin{proof}\leanok
    Let $k=\lvert\{z\in Z:p_z\neq0\}\rvert$. Jensen's inequality for the
    concave function $h$, with uniform weights on the support, gives
    \begin{align}
        \frac{1}{k}\sum_{\substack{z\in Z \\p_z\neq0}}h(p_z)
         & \leq
        h\left(\frac{1}{k}
        \sum_{\substack{z\in Z            \\p_z\neq0}}p_z\right)
        =h\left(\frac{1}{k}\right)
        =\frac{\log k}{k}.
        \notag
    \end{align}
    Multiplication by $k$ proves the result, since values outside the support
    contribute $h(0)=0$.
\end{proof}

\begin{theorem}[Classical mutual information is bounded by ordinary rank]
    \label{thm:classical_mutual_information_le_log_rank}
    \lean{Entropy.classicalMutualInformation_le_log_rank}
    \leanok
    \uses{def:finite_classical_mutual_information}
    Let $X$ and $Y$ be finite sets and let $P=(P_{x,y})_{x\in X,y\in Y}$
    be a joint probability distribution. With natural logarithms,
    $I(X:Y)_P\leq\log\rank_{\R}P$. Equivalently, with logarithms to base two,
    $2^{I(X:Y)_P}\leq\rank_{\R}P$.
\end{theorem}
\begin{proof}\leanok
    \uses{lem:probability_entropy_le_log_card_support}
    This is Theorem~4.1 of \cite{Riazanov2017PSDRank}. If the number of
    nonzero rows exceeds the ordinary real rank, choose a nontrivial real
    linear relation among those rows. Nonnegativity forces the relation to
    have coefficients of both signs. Rescaling each row by
    $1+\varepsilon\beta_x$ gives two endpoint distributions. For every
    $y\in Y$, the row relation gives
    \begin{align}
        \sum_x(1+\varepsilon\beta_x)P_{x,y}
         & =\sum_xP_{x,y}
        +\varepsilon\underbrace{\sum_x\beta_xP_{x,y}}_{=0}
        =\sum_xP_{x,y}.
        \notag
    \end{align}
    Thus each endpoint preserves every column marginal, removes at least one
    row, and has no larger rank.
    Write $h(t)=-t\log t$ for $t>0$, with $h(0)=0$, and put
    $m=\min_x\beta_x<0<M=\max_x\beta_x$. Then
    \begin{align}
        I(X:Y)_{P^{(\varepsilon)}}
         & =I(X:Y)_P+\varepsilon S(\beta),
        \notag                                                  \\
        S(\beta)
         & =\sum_x\beta_x\left(h(p_x)-\sum_y h(P_{x,y})\right).
        \notag
    \end{align}
    If $S(\beta)\geq0$, choose $\varepsilon=-m^{-1}\geq0$; otherwise choose
    $\varepsilon=-M^{-1}\leq0$. In either case
    $\varepsilon S(\beta)\geq0$, so the selected endpoint has mutual
    information no smaller than the original distribution.
    Repeating the construction leaves at most
    $\rank_{\R}P$ nonzero rows. Finally,
    $I(X:Y)\leq H(X)$, and the entropy of a distribution supported on $k$
    points is at most $\log k$.
\end{proof}
\subsection{Operator-Schmidt rank and marginal-support compression}

\begin{theorem}[Positivity of the operator-Schmidt rank]
    \label{thm:operator_schmidt_rank_pos}
    \lean{Matrix.operatorSchmidtRank_pos_of_ne_zero}
    \leanok

    Every nonzero finite-dimensional bipartite operator $X$ satisfies
    \begin{align}
        0 & <\operatorname{OSR}(X).
        \notag
    \end{align}
    This includes zero-dimensional factors, where the nonzero hypothesis is
    impossible.
\end{theorem}

\begin{proof}\leanok

    If $\operatorname{OSR}(X)=0$, a shortest product decomposition of $X$ is
    an empty sum. Hence $X=0$.
\end{proof}
\begin{theorem}[Coefficient-space characterization of operator-Schmidt rank]
    \label{thm:operator_schmidt_coefficient_space}
    \lean{Matrix.range_operatorSchmidtMap_eq_operatorBlockSpan,
        Matrix.operatorSchmidtRank_eq_finrank_operatorBlockSpan}
    \leanok

    Write $\rho_{ij}\in\MN{d_B}$ for the blocks determined by a basis
    of the first factor, and define
    $\mathcal R_\rho(X)=\sum_{i,j}X_{ij}\rho_{ij}$. Then
    \begin{align}
        \operatorname{OSR}(\rho)
         & =\dim\range\mathcal R_\rho
        =\dim\spn\{\rho_{ij}:1\leq i,j\leq d_A\}.
        \notag
    \end{align}
\end{theorem}
\begin{proof}\leanok

    The minimality theorem at the preceding definition proves that the range
    dimension $\dim\range\mathcal R_\rho$ is the least admissible
    product-decomposition length: it is itself an admissible length and no
    shorter length suffices. Since the operator-Schmidt rank is defined as the
    least such length, uniqueness of the least integer identifies the two, giving
    the first equality $\operatorname{OSR}(\rho)=\dim\range\mathcal R_\rho$. For
    the second equality, $\mathcal R_\rho$ is the linear combination map sending
    a coefficient matrix $X$ to $\sum_{i,j}X_{ij}\rho_{ij}$, so its range is
    exactly the span of the operator blocks $\rho_{ij}$.
\end{proof}

\begin{theorem}[Operator-Schmidt rank under local isometries]
    \label{thm:operator_schmidt_rank_local_isometry}
    \lean{Matrix.operatorSchmidtRank_local_mul_le,
        Matrix.operatorSchmidtRank_local_isometry_conj}
    \leanok

    Let $V_A:\mathcal H_A\to\mathcal K_A$ and
    $V_B:\mathcal H_B\to\mathcal K_B$ be isometries between finite-dimensional
    complex spaces. For every operator $X$ on
    $\mathcal H_A\otimes\mathcal H_B$,
    \begin{align}
        \operatorname{OSR}\!\left(
        (V_A\otimes V_B)X(V_A\otimes V_B)^\dagger\right)
         & =\operatorname{OSR}(X).
        \notag
    \end{align}
    This remains valid when one or more of the spaces have dimension zero.
\end{theorem}
\begin{proof}\leanok
    \uses{thm:operator_schmidt_coefficient_space}
    More generally, multiplying on the left and right by product matrices
    cannot increase operator-Schmidt rank: apply the four local matrices to
    the two factors in each term of a shortest product decomposition. Applying
    this inequality first to $V_A,V_B$ and then to their adjoints gives the two
    inequalities. The identities $V_A^\dagger V_A=1$ and
    $V_B^\dagger V_B=1$ identify the twice-transformed operator with $X$.
\end{proof}

\begin{theorem}[Support of a bipartite positive operator]
    \label{thm:product_marginals_kernel}
    \lean{Matrix.PosSemidef.productMarginals_kernel_le}
    \leanok
    Let $\rho_{AB}\succeq0$, with marginals $\rho_A$ and $\rho_B$. Then
    \begin{align}
        \ker(\rho_A\otimes\rho_B) & \subseteq\ker\rho_{AB}.
        \notag
    \end{align}
    Equivalently, the support of $\rho_{AB}$ is contained in
    $\supp(\rho_A)\otimes\supp(\rho_B)$. This remains valid when either
    factor has dimension zero.
\end{theorem}

\begin{proof}\leanok
    Let $P_A$ and $P_B$ be the support projections of the marginals. The two
    right-absorption identities give
    $\rho_{AB}(P_A\otimes P_B)=\rho_{AB}$. The support projection of
    $\rho_A\otimes\rho_B$ is $P_A\otimes P_B$. Therefore every vector
    annihilated by $\rho_A\otimes\rho_B$ is annihilated by $\rho_{AB}$.
\end{proof}

\begin{theorem}[Compression to both marginal supports]
    \label{thm:operator_schmidt_rank_marginal_support_compression}
    \lean{Matrix.PosSemidef.eq_marginalSupport_expansion_compression,
        Matrix.PosSemidef.operatorSchmidtRank_marginalSupport_compression}
    \leanok

    Let $\rho\geq0$ be an operator on $H_A\otimes H_B$. Let
    $V_A:\widehat H_A\to H_A$ and $V_B:\widehat H_B\to H_B$ be isometries
    whose range projectors are the support projectors $P_A$ and $P_B$ of the
    two marginals. Set $W=V_A\otimes V_B$ and
    $\rho_c=W^\dagger\rho W$. Then
    \begin{align}
        W\rho_cW^\dagger & =\rho,
                         & \operatorname{OSR}(\rho_c) & =\operatorname{OSR}(\rho).
        \notag
    \end{align}
    No marginal is required to be faithful, and zero-dimensional coordinate
    spaces are permitted.
\end{theorem}
\begin{proof}\leanok
    \uses{thm:operator_schmidt_rank_local_isometry}
    The four marginal-support absorption identities give
    \begin{align}
        (P_A\otimes P_B)\rho & =\rho,
                             & \rho(P_A\otimes P_B) & =\rho.
        \notag
    \end{align}
    Since $WW^\dagger=P_A\otimes P_B$, associativity yields
    \begin{align}
        W\rho_cW^\dagger
         & =(WW^\dagger)\rho(WW^\dagger)=\rho.
        \notag
    \end{align}
    The local-isometry theorem applied to $\rho_c$ gives
    $\operatorname{OSR}(W\rho_cW^\dagger)=\operatorname{OSR}(\rho_c)$.
    Substitution of the reconstruction identity proves the rank equality.
\end{proof}

\begin{lemma}[Partial traces under product isometries]
    \label{lem:partial_trace_product_isometry}
    \lean{Matrix.partialTraceRight_kronecker_conj_of_right_isometry,
        Matrix.partialTraceLeft_kronecker_conj_of_left_isometry}
    \leanok
    Let $A:H_A\to K_A$ and $B:H_B\to K_B$ be linear maps between
    finite-dimensional complex spaces, and let $X$ be an operator on
    $H_A\otimes H_B$. If $B^\dagger B=1$, then
    \begin{align}
        \tr_{K_B}\!\left((A\otimes B)X(A\otimes B)^\dagger\right)
         & =A(\tr_{H_B}X)A^\dagger.
        \notag
    \end{align}
    If $A^\dagger A=1$, then, symmetrically,
    \begin{align}
        \tr_{K_A}\!\left((A\otimes B)X(A\otimes B)^\dagger\right)
         & =B(\tr_{H_A}X)B^\dagger.
        \notag
    \end{align}
    The identities include zero-dimensional spaces.
\end{lemma}
\begin{proof}\leanok
    Expand the matrix entries in orthonormal bases. In the first identity, the
    sum over the traced output basis gives the matrix entries of
    $B^\dagger B=1$, and hence collapses the two input indices on $H_B$.
    The second identity follows by the same argument with the factors
    interchanged.
\end{proof}

\begin{theorem}[Faithful marginals after simultaneous support compression]
    \label{thm:faithful_marginals_support_compression}
    \lean{Matrix.PosSemidef.partialTraceRight_marginalSupport_compression,
        Matrix.PosSemidef.partialTraceLeft_marginalSupport_compression,
        Matrix.PosSemidef.marginalSupport_compression_marginals_posDef}
    \leanok
    \uses{thm:operator_schmidt_rank_marginal_support_compression}
    In the setting of
    Theorem~\ref{thm:operator_schmidt_rank_marginal_support_compression}, the
    two marginals of $\rho_c$ satisfy
    \begin{align}
        \tr_{\widehat H_B}\rho_c
         & =V_A^\dagger(\tr_{H_B}\rho)V_A,
         &
        \tr_{\widehat H_A}\rho_c
         & =V_B^\dagger(\tr_{H_A}\rho)V_B.
        \notag
    \end{align}
    Both compressed marginals are positive definite, including when one of
    their spaces has dimension zero.
\end{theorem}
\begin{proof}\leanok
    \uses{lem:partial_trace_product_isometry, thm:operator_schmidt_rank_marginal_support_compression}
    Insert the reconstruction $\rho=W\rho_cW^\dagger$ into the two covariance
    identities of Lemma~\ref{lem:partial_trace_product_isometry}, and multiply
    by the corresponding adjoint isometries. Each displayed marginal is the
    compression of a positive semidefinite matrix to its support, so its
    positive definiteness follows from
    Theorem~\ref{thm:compression_on_support_posDef}.
\end{proof}

\begin{theorem}[Channel of a faithful bipartite marginal]
    \label{thm:supported_marginal_channel}
    \lean{Matrix.finrank_range_supportedMarginalMap,
        Matrix.supportedMarginalMap_isKrausCPTP,
        Matrix.supportedMarginalMap_diagonal,
        Matrix.supportedMarginalReconstruction_eq}
    \leanok
    \uses{thm:operator_schmidt_coefficient_space}
    Let $\rho\geq0$ be a bipartite complex matrix whose first marginal is
    faithful. Choose an eigenbasis in which
    $\tr_B\rho=\diag(p_1,\ldots,p_{d_A})$ with
    every $p_i>0$, and define the linear map $\Phi_\rho$ on matrix units by
    \begin{align}
        \Phi_\rho(E_{ij})
         & =\frac{\rho_{ij}}{\sqrt{p_ip_j}}.
        \notag
    \end{align}
    Write $\sigma=\diag(p_1,\ldots,p_{d_A})$ and
    $\tau=\tr_A\rho$. Then $\Phi_\rho$ is completely positive and trace
    preserving,
    \begin{align}
        \Phi_\rho(\sigma) & =\tau,
                          & \dim\range\Phi_\rho & =\operatorname{OSR}(\rho),
        \notag
    \end{align}
    and application of
    $\Phi_\rho$ to the first half of
    $\sum_{i,j}\sqrt{p_ip_j}\,E_{ij}\otimes E_{ij}$, followed by restoring the
    order of the two factors, reconstructs $\rho$.
\end{theorem}
\begin{proof}\leanok
    Strict positivity of the $p_i$ makes the entrywise input scaling
    invertible, so it does not change the range. Positivity of $\rho$ gives a
    Kraus representation of $\Phi_\rho$, while the first marginal equation
    gives trace preservation. Summing the diagonal blocks gives
    $\Phi_\rho(\sigma)=\tau$. Substitution on matrix units proves the
    reconstruction identity. This is the finite-dimensional Choi
    representation \cite{Choi1975CompletelyPositive} with the input marginal
    absorbed into the canonical purification.
\end{proof}

\section{Support compression for entropy functionals}

% Issues #5229, #5241, and #5245 establish one- and two-marginal support
% compression, including preservation of operator-Schmidt rank and
% faithfulness of both compressed marginals.  Issue #5211 combines these
% results with the whitened-Choi estimate; #5212 tracks mutual information.

\begin{definition}[Isometric conjugation into a larger matrix algebra]
    \label{def:isometric_nonunital_conjugation}
    \lean{Matrix.isometryConjNonUnitalStarAlgHom,
        Matrix.isometryConjNonUnitalStarAlgHom_apply}
    \leanok
    Let $V:\C^k\to\C^D$ be an isometry, so that
    $V^\dagger V=\Id_k$. Define $\iota_V:\MN{k}\to\MN{D}$ by
    \begin{align}
        \iota_V(A) & =VAV^\dagger.
        \label{eq:mpdo_isometric_conjugation}
    \end{align}
    The map $\iota_V$ is complex-linear, multiplicative, and $*$-preserving.
    In general it is not unital:
    $\iota_V(\Id_k)=VV^\dagger$ is the projection onto the range of $V$.
\end{definition}

\begin{theorem}[Functional calculus under rectangular isometric conjugation]
    \label{thm:cfc_rectangular_isometry_zero}
    \lean{Matrix.cfc_conj_isometry_of_zero}
    \leanok
    Let $A\in\MN{k}$ be Hermitian, let $V:\C^k\to\C^D$ satisfy
    $V^\dagger V=\Id_k$, and let $f:\R\to\R$ satisfy $f(0)=0$. Then
    \begin{align}
        f(VAV^\dagger) & =Vf(A)V^\dagger,
        \label{eq:mpdo_cfc_isometric_conjugation}
    \end{align}
    where both sides use the continuous functional calculus on the respective
    finite spectra.
\end{theorem}

\begin{proof}\leanok
    \uses{def:isometric_nonunital_conjugation}
    Apply functoriality of the continuous functional calculus to
    $\iota_V$. The condition $f(0)=0$ removes the contribution from the
    orthogonal complement of the range of $V$, where $\iota_V(A)$ vanishes.
\end{proof}

\begin{theorem}[Trace preservation under isometric inclusion]
    \label{thm:isometric_embedding_trace}
    \lean{Matrix.trace_mul_mul_conjTranspose_of_conjTranspose_mul_eq_one}
    \leanok
    Let $V:\C^k\to\C^D$ satisfy $V^\dagger V=\Id_k$. For every
    $A\in\MN{k}$,
    \begin{align}
        \tr(VAV^\dagger) & =\tr(A).
        \label{eq:mpdo_isometric_trace}
    \end{align}
\end{theorem}

\begin{proof}\leanok
    Cyclicity of the trace gives
    $\tr(VAV^\dagger)=\tr(V^\dagger VA)=\tr(A)$.
\end{proof}

\begin{theorem}[Compression cancels isometric expansion]
    \label{thm:isometric_compression_expansion}
    \lean{Matrix.conjTranspose_mul_mul_mul_conjTranspose_mul_of_isometry}
    \leanok
    Let $V:\C^k\to\C^D$ satisfy $V^\dagger V=\Id_k$. For every
    $A\in\MN{k}$,
    \begin{align}
        V^\dagger(VAV^\dagger)V & =A.
        \label{eq:mpdo_isometric_compression_expansion}
    \end{align}
\end{theorem}

\begin{proof}\leanok
    Associativity and $V^\dagger V=\Id_k$ reduce the left-hand side to
    $(V^\dagger V)A(V^\dagger V)=A$.
\end{proof}

\begin{theorem}[Reconstruction from a reference support]
    \label{thm:support_compression_reconstruction}
    \lean{Matrix.PosSemidef.eq_expansion_compression_of_kernel_le_of_mul_conjTranspose_eq_supportProj}
    \leanok
    Let $\rho,\omega\in\MN{D}$ be positive semidefinite and suppose that
    $\ker\omega\subseteq\ker\rho$. If
    $V:\C^k\to\C^D$ satisfies $VV^\dagger=P_{\supp(\omega)}$, then
    \begin{align}
        V(V^\dagger\rho V)V^\dagger & =\rho.
        \label{eq:mpdo_support_reconstruction}
    \end{align}
    No condition on $V^\dagger V$ is needed for this identity.
\end{theorem}

\begin{proof}\leanok
    The kernel inclusion and positivity give
    $\rho P_{\supp(\omega)}=\rho$; taking adjoints also gives
    $P_{\supp(\omega)}\rho=\rho$. Substituting
    $VV^\dagger=P_{\supp(\omega)}$ on both sides of $\rho$ proves
    (\ref{eq:mpdo_support_reconstruction}).
\end{proof}

\begin{theorem}[Relative entropy under support compression]
    \label{thm:relative_entropy_support_compression}
    \lean{TNLean.quantumRelativeEntropy_support_compression}
    \leanok
    Let $\rho,\omega\in\MN{D}$ be positive semidefinite with
    $\ker\omega\subseteq\ker\rho$. Let $V:\C^k\to\C^D$ satisfy
    \begin{align}
        V^\dagger V & =\Id_k,             \\
        VV^\dagger  & =P_{\supp(\omega)}.
        \label{eq:mpdo_reference_support_isometry}
    \end{align}
    Then
    \begin{align}
        D(\rho\Vert\omega)
         & =D(V^\dagger\rho V\Vert V^\dagger\omega V).
        \label{eq:mpdo_relative_entropy_support_compression}
    \end{align}
    The logarithm is totalized by $\log 0=0$ on the kernel.
\end{theorem}

\begin{proof}\leanok
    \uses{def:isometric_nonunital_conjugation, thm:cfc_rectangular_isometry_zero, thm:isometric_embedding_trace, thm:support_compression_reconstruction}
    Theorem~\ref{thm:support_compression_reconstruction} reconstructs both
    $\rho$ and $\omega$ from their compressions. Apply
    Theorem~\ref{thm:cfc_rectangular_isometry_zero} to the logarithm, use
    multiplicativity of $\iota_V$, and finish with
    (\ref{eq:mpdo_isometric_trace}).
\end{proof}

\begin{theorem}[Functional calculus under transpose]
    \label{thm:cfc_transpose_hermitian}
    \lean{Matrix.cfc_transpose}
    \lean{Matrix.IsHermitian.log_transpose}
    \lean{Matrix.PosSemidef.rpow_transpose}
    \lean{Matrix.PosSemidef.sqrt_transpose}
    \lean{Matrix.IsHermitian.supportProj_transpose}
    \leanok
    If $A\in\MN{n}$ is Hermitian and $f:\R\to\R$, then
    \begin{align}
        f(A)^T & =f(A^T).
        \label{eq:mpdo_cfc_transpose}
    \end{align}
    In particular,
    \begin{align}
        \log(A^T) & =(\log A)^T, &
        P_{A^T}   & =P_A^T.
        \label{eq:mpdo_log_support_transpose}
    \end{align}
    If $A$ is positive semidefinite, then for every $r\in\R$,
    \begin{align}
        (A^T)^r    & =(A^r)^T,     &
        \sqrt{A^T} & =(\sqrt A)^T.
        \label{eq:mpdo_power_sqrt_transpose}
    \end{align}
    These identities include singular matrices and matrices whose index set is
    empty. They are project-derived rather than statements of CPSV16.
\end{theorem}

\begin{proof}\leanok
    For a Hermitian matrix, transpose is entrywise complex conjugation.
    Apply covariance of the continuous functional calculus under this real
    star-algebra automorphism. Continuity of $f$ is needed only on the finite
    spectrum of $A$, where it is automatic. The logarithm, real-power, and
    square-root identities follow by specialization. For the support
    projection, specialize to $f(x)=1$ for $x\neq0$ and $f(0)=0$.
\end{proof}


\begin{theorem}[Logarithm of a positive square root]
    \label{thm:log_sqrt_possemidef}
    \lean{Matrix.PosSemidef.cfc_log_sqrt}
    \leanok
    If $\rho\succeq0$, then
    \begin{align}
        \log\sqrt{\rho} & =\frac12\log\rho.
        \label{eq:mpdo_log_sqrt}
    \end{align}
\end{theorem}

\begin{proof}\leanok
    On every positive eigenvalue this is the scalar identity
    $\log\sqrt{x}=\frac12\log x$. At $x=0$, both sides vanish under the
    convention $\log0=0$.
\end{proof}

\begin{theorem}[Logarithm of a real power]
    \label{thm:log_rpow_posdef}
    \lean{Matrix.PosDef.cfc_log_rpow}
    \leanok
    If $A\succ0$, then, for every $s\in\R$,
    \begin{align}
        \log(A^s) & =s\log A.
        \label{eq:mpdo_log_rpow}
    \end{align}
\end{theorem}

\begin{proof}\leanok
    Every eigenvalue of $A$ is positive, so the assertion follows from
    $\log(x^s)=s\log x$ on $(0,\infty)$ and the continuous functional
    calculus.
\end{proof}

\begin{theorem}[Support-aware vector-state Jensen inequality for the logarithm]
    \label{thm:support_log_vector_jensen}
    \lean{Matrix.PosSemidef.re_dotProduct_cfc_log_mulVec_le_log}
    \leanok
    Let $A\in\MN{n}$ be positive semidefinite and let $v\in\C^n$ satisfy
    $\langle v,v\rangle=1$ and $P_{\supp(A)}v=v$. Then
    \begin{align}
        \re\langle v,(\log A)v\rangle
         & \leq\log\!\left(\re\langle v,Av\rangle\right).
        \label{eq:mpdo_support_log_jensen}
    \end{align}
    Here $\log A$ is defined by the continuous functional calculus with
    $\log 0=0$. Compare the finite-dimensional vector-state Jensen argument
    in \cite[Chapter~V]{Bhatia1997Matrix}.
\end{theorem}

\begin{proof}\leanok
    Diagonalize $A=U\diag(\mu_i)U^\dagger$, set $w=U^\dagger v$, and put
    $p_i=|w_i|^2$. The support condition implies
    $\mu_i=0\Longrightarrow w_i=0$, so every spectral weight at zero vanishes.
    Thus $p_i\geq0$ for every $i\in S=\{i:\mu_i>0\}$ and
    $\sum_{i\in S}p_i=1$. The spectral formulas give
    \begin{align}
        \re\langle v,(\log A)v\rangle
         & =\sum_{i\in S}p_i\log\mu_i, \notag \\
        \re\langle v,Av\rangle
         & =\sum_{i\in S}p_i\mu_i.
        \notag
    \end{align}
    Scalar concavity of the logarithm on $(0,\infty)$ now gives
    (\ref{eq:mpdo_support_log_jensen}). No concavity assertion at zero and no
    positive-definiteness of $A$ are used.
\end{proof}

\begin{definition}[Order-two sandwiched trace functional]
    \label{def:sandwiched_renyi_two_trace}
    \lean{TNLean.sandwichedRenyiTwoTrace}
    \leanok
    For square matrices $\rho$ and $\omega$ of the same size, define
    \begin{align}
        Q_2(\rho,\omega)
         & =\re\tr\left[
            \left(\omega^{-1/4}\rho\,\omega^{-1/4}\right)^2
            \right].
        \label{eq:mpdo_sandwiched_two_trace}
    \end{align}
    The negative power is defined by functional calculus and vanishes on the
    kernel of $\omega$. On trace-one positive semidefinite matrices satisfying
    $\ker\omega\subseteq\ker\rho$, its logarithm is the order-two sandwiched
    R\'enyi divergence \cite[Definition~2]{MullerLennert2013Renyi}.
\end{definition}

\begin{theorem}[Nonnegativity of the order-two functional]
    \label{thm:sandwiched_renyi_two_nonneg}
    \lean{TNLean.sandwichedRenyiTwoTrace_nonneg}
    \leanok
    \uses{def:sandwiched_renyi_two_trace}
    If $\rho\succeq0$ and $\omega\succeq0$, then
    $Q_2(\rho,\omega)\geq0$.
\end{theorem}

\begin{proof}\leanok
    \uses{def:sandwiched_renyi_two_trace}
    The sandwiched matrix
    $\omega^{-1/4}\rho\,\omega^{-1/4}$ is positive semidefinite. The trace of
    its square is therefore nonnegative.
\end{proof}

\begin{theorem}[Faithful cyclic form of the order-two functional]
    \label{thm:sandwiched_renyi_two_weighted}
    \lean{TNLean.sandwichedRenyiTwoTrace_eq_weighted}
    \leanok
    \uses{def:sandwiched_renyi_two_trace}
    If $\omega\succ0$, then, for every square matrix $\rho$ of the same size,
    \begin{align}
        Q_2(\rho,\omega)
         & =\re\tr\left(\rho\,\omega^{-1/2}
        \rho\,\omega^{-1/2}\right).
        \label{eq:mpdo_sandwiched_two_weighted}
    \end{align}
\end{theorem}

\begin{proof}\leanok
    Combine the two adjacent inverse quarter-powers by
    $\omega^{-1/4}\omega^{-1/4}=\omega^{-1/2}$ and rotate the factors under
    the trace. No Hermiticity assumption on $\rho$ is needed.
\end{proof}

\begin{theorem}[Relative entropy bounded by the faithful order-two functional]
    \label{thm:relative_entropy_q2_posdef}
    \lean{TNLean.quantumRelativeEntropy_le_log_sandwichedRenyiTwoTrace_posDef}
    \leanok
    \uses{def:sandwiched_renyi_two_trace}
    Let $\rho,\omega\in\MN{D}$ satisfy $\rho\succeq0$, $\omega\succ0$, and
    $\tr\rho=1$. Then
    \begin{align}
        D(\rho\Vert\omega) & \leq\log Q_2(\rho,\omega).
        \label{eq:mpdo_relative_q2_posdef}
    \end{align}
    No normalization of $\omega$ is required.
\end{theorem}

\begin{proof}\leanok
    \uses{thm:cfc_transpose_hermitian, thm:log_sqrt_possemidef, thm:log_rpow_posdef, thm:support_log_vector_jensen, thm:sandwiched_renyi_two_weighted}
    Set $R=\sqrt\rho$, $T=\omega^{-1/2}$,
    $\Delta=R^T\otimes T$, and $v=\operatorname{vec}(R)$. The vectorization
    is column-stacking, with factor order
    \begin{align}
        (B\otimes A)\operatorname{vec}(X)
         & =\operatorname{vec}(AXB^T).
        \label{eq:mpdo_column_vec_kronecker}
    \end{align}
    Hence $\Delta v=\operatorname{vec}(T\rho)$, and cyclicity of the trace
    identifies its first moment with
    \begin{align}
        m=\re\langle v,\Delta v\rangle
         & =\langle R,RTR\rangle_{\mathrm{HS}}.
        \label{eq:mpdo_relative_modular_half_moment}
    \end{align}
    The support projection of $\Delta$ fixes $v$:
    \begin{align}
        P_{\supp(\Delta)}v
         & =(P_{\supp(R^T)}\otimes P_{\supp(T)})\operatorname{vec}(R)
        =\operatorname{vec}(R)=v.
        \label{eq:mpdo_relative_modular_support_fixing}
    \end{align}
    Indeed, $T$ is positive definite, so $P_{\supp(T)}=1$, while the support
    projection of $R^T$ fixes $R^T$ and hence
    $R P_{\supp(R^T)}^T=R$. The column-stacking identity
    (\ref{eq:mpdo_column_vec_kronecker}) proves the displayed equality. This
    support condition is essential when $\rho$ is singular: it removes the
    zero spectral weights before Jensen's inequality is applied.

    Theorem~\ref{thm:mpdo_support_log_kronecker},
    \ref{thm:log_sqrt_possemidef}, \ref{thm:log_rpow_posdef}, and
    \ref{thm:cfc_transpose_hermitian} give
    \begin{align}
        \re\langle v,(\log\Delta)v\rangle
         & =\frac12 D(\rho\Vert\omega).
        \label{eq:mpdo_relative_modular_log_moment}
    \end{align}
    Theorem~\ref{thm:support_log_vector_jensen} therefore yields
    $\frac12D(\rho\Vert\omega)\leq\log m$. Since
    $\|R\|_{\mathrm{HS}}^2=\tr\rho=1$, Hilbert--Schmidt
    Cauchy--Schwarz gives
    \begin{align}
        m^2
         & \leq\|R\|_{\mathrm{HS}}^2\|RTR\|_{\mathrm{HS}}^2
        =Q_2(\rho,\omega).
        \label{eq:mpdo_relative_modular_cauchy_schwarz}
    \end{align}
    Positivity of $m$ now gives
    $D(\rho\Vert\omega)\leq\log(m^2)\leq\log Q_2(\rho,\omega)$.

    The auxiliary divergence and its vector-state Jensen argument are adapted
    from \cite[arXiv:1306.3142v4, Definition~5 and Lemma~19, used in the proof
        of Theorem~7]{MullerLennert2013Renyi}.  Only the direct order-two endpoint
    is used here; the final Cauchy--Schwarz estimate is not an application of
    Beigi's interpolation theorem.
\end{proof}

\begin{theorem}[Order-two sandwiched functional under support compression]
    \label{thm:sandwiched_renyi_two_support_compression}
    \lean{TNLean.sandwichedRenyiTwoTrace_support_compression}
    \leanok
    \uses{def:sandwiched_renyi_two_trace}
    Let $\rho,\omega\in\MN{D}$ be positive semidefinite with
    $\ker\omega\subseteq\ker\rho$. Let $V:\C^k\to\C^D$ satisfy
    (\ref{eq:mpdo_reference_support_isometry}). Then
    \begin{align}
        Q_2(\rho,\omega)
         & =Q_2(V^\dagger\rho V,V^\dagger\omega V).
        \label{eq:mpdo_sandwiched_two_support_compression}
    \end{align}
\end{theorem}

\begin{proof}\leanok
    \uses{def:isometric_nonunital_conjugation, thm:cfc_rectangular_isometry_zero, thm:isometric_embedding_trace, thm:support_compression_reconstruction}
    The function $x\mapsto x^{-1/4}$ is taken to be zero at $x=0$, so
    Theorem~\ref{thm:cfc_rectangular_isometry_zero} applies. Reconstruct
    $\rho$ and $\omega$, combine the three expanded factors by
    multiplicativity of $\iota_V$, and use trace preservation.
\end{proof}

\begin{theorem}[Reduction of the singular order-two comparison to the faithful case]
    \label{thm:relative_entropy_q2_faithful_reduction}
    \lean{TNLean.quantumRelativeEntropy_le_log_sandwichedRenyiTwoTrace_of_faithful}
    \leanok
    \uses{def:sandwiched_renyi_two_trace}
    Let $\rho,\omega\in\MN{D}$ be positive semidefinite matrices of trace one
    with $\ker\omega\subseteq\ker\rho$. Assume that for every dimension and
    every pair of trace-one matrices $A\succeq0$ and $B\succ0$,
    \begin{align}
        D(A\Vert B) & \leq\log Q_2(A,B).
        \label{eq:mpdo_faithful_relative_q2_hypothesis}
    \end{align}
    Then
    \begin{align}
        D(\rho\Vert\omega) & \leq\log Q_2(\rho,\omega).
        \label{eq:mpdo_singular_relative_q2_comparison}
    \end{align}
    Thus this theorem reduces the singular-reference case to the faithful
    comparison. Theorem~\ref{thm:relative_entropy_q2_posdef} supplies the
    hypothesis in~(\ref{eq:mpdo_faithful_relative_q2_hypothesis}).
\end{theorem}

\begin{proof}\leanok
    \uses{thm:isometric_embedding_trace, thm:support_compression_reconstruction, thm:relative_entropy_support_compression, thm:sandwiched_renyi_two_support_compression}
    Choose an isometric inclusion of the support of $\omega$. Its compression
    $\omega_V=V^\dagger\omega V$ is positive definite by
    Theorem~\ref{thm:compression_on_support_posDef}, while
    $\rho_V=V^\dagger\rho V$ remains positive semidefinite. Reconstruction
    and (\ref{eq:mpdo_isometric_trace}) show that both compressed matrices
    have trace one. Apply (\ref{eq:mpdo_faithful_relative_q2_hypothesis}) to
    $(\rho_V,\omega_V)$, then use
    (\ref{eq:mpdo_relative_entropy_support_compression}) and
    (\ref{eq:mpdo_sandwiched_two_support_compression}).
\end{proof}

\begin{theorem}[Relative entropy bounded by the support-restricted order-two functional]
    \label{thm:relative_entropy_q2_support}
    \lean{TNLean.quantumRelativeEntropy_le_log_sandwichedRenyiTwoTrace}
    \leanok
    \uses{def:sandwiched_renyi_two_trace}
    Let $\rho,\omega\in\MN{D}$ be positive semidefinite matrices of trace one.
    If $\ker\omega\subseteq\ker\rho$, then
    \begin{align}
        D(\rho\Vert\omega) & \leq\log Q_2(\rho,\omega).
        \label{eq:mpdo_relative_q2_support}
    \end{align}
\end{theorem}

\begin{proof}\leanok
    \uses{thm:relative_entropy_q2_posdef, thm:relative_entropy_q2_faithful_reduction}
    Apply Theorem~\ref{thm:relative_entropy_q2_faithful_reduction} with the
    faithful comparison from Theorem~\ref{thm:relative_entropy_q2_posdef}.
    Compression to the support of $\omega$ preserves both entropy functionals
    and both trace normalizations, and makes the reference positive definite.
\end{proof}
\subsection{Whitened Choi estimates and sandwiched R\'enyi bounds}

\begin{theorem}[Full-support eigenbasis whitening as a Choi matrix]
    \label{thm:full_support_eigenbasis_whitening_choi}
    \lean{Matrix.rectangularChoi_singleKrausMap_comp_comp,
        Matrix.supportedMarginalWhitenedState_eq_rectangularChoi,
        Matrix.finrank_range_weightedHilbertSchmidtMap}
    \leanok
    \uses{thm:supported_marginal_channel, def:rectangular_choi}
    Let $\rho_{AB}\geq0$, and choose an eigenbasis of its faithful first
    marginal
    $\sigma=\diag(p_1,\ldots,p_{d_A})>0$.  Suppose also that the second
    marginal $\tau$ is positive definite.  For the supported-marginal channel
    $\Phi_\rho$, set
    \begin{align}
        L(X) & =\tau^{-1/4}
        \Phi_\rho(\sigma^{1/4}X\sigma^{1/4})\tau^{-1/4}, \\
        W    & =(\sigma^T\otimes\tau)^{-1/4}\rho_{AB}
        (\sigma^T\otimes\tau)^{-1/4}.
        \notag
    \end{align}
    Then
    \begin{align}
        W & =J(L),
          & \dim_{\C}\range L & =\dim_{\C}\range\Phi_\rho.
        \notag
    \end{align}
    Here $\sigma^T=\sigma$ because the chosen basis diagonalizes $\sigma$.
\end{theorem}

\begin{proof}\leanok
    \uses{thm:supported_marginal_channel}
    For arbitrary congruence matrices $A$ and $B$, the matrix-unit convention
    for the Choi matrix gives
    \begin{align}
        J(\Gamma_B\circ\Phi\circ\Gamma_A)
         & =\Gamma_{A^T\otimes B}(J(\Phi)).
        \notag
    \end{align}
    Apply this identity to the inverse-square-root scaling in
    $\Phi_\rho$ and to the two fourth-root factors defining $L$.
    Entrywise diagonal functional calculus gives
    $\sigma^{1/4}=\diag(p_i^{1/4})$ and
    $\diag(p_i^{-1/2})\sigma^{1/4}=\sigma^{-1/4}$, proving the first
    assertion.  The two modular congruences are invertible linear maps, so
    pre- and post-composition by them preserve the range dimension.
\end{proof}

\begin{theorem}[Full-support eigenbasis whitened Choi estimate]
    \label{thm:full_support_eigenbasis_whitened_choi_bound}
    \lean{Matrix.supportedMarginalWhitenedState_posSemidef,
        Matrix.supportedMarginalWhitenedState_frobenius_sq_le_operatorSchmidtRank,
        Matrix.supportedMarginalWhitenedState_trace_sq_re_le_operatorSchmidtRank}
    \leanok
    \uses{thm:full_support_eigenbasis_whitening_choi, thm:rectangular_choi_rank_bound}
    Under the hypotheses of
    Theorem~\ref{thm:full_support_eigenbasis_whitening_choi},
    the matrix
    \begin{align}
        W=(\sigma^T\otimes\tau)^{-1/4}\rho_{AB}
        (\sigma^T\otimes\tau)^{-1/4}
        \notag
    \end{align}
    is positive semidefinite, and
    \begin{align}
        \tr(W^2)=\|W\|_F^2
         & \leq \operatorname{OSR}(\rho_{AB}).
        \notag
    \end{align}
    This is the faithful-marginal-support form of the order-two estimate
    obtained from \cite[Theorem~6, Equation~(18)]{Beigi2013Sandwiched}.
    % The passage from singular ambient marginals to their supports remains
    % recorded in docs/paper-gaps/cpgsv17_mpdo_mutual_information_bound.tex.
\end{theorem}

\begin{proof}\leanok
    \uses{thm:full_support_eigenbasis_whitening_choi, thm:rectangular_choi_rank_bound}
    The matrix $W$ is a congruence of $\rho_{AB}\geq0$, so it is positive
    semidefinite and Hermitian; consequently
    $\tr(W^2)=\|W\|_F^2$.
    The weighted map $L$ is a Hilbert--Schmidt contraction.  Its Choi matrix
    is the whitened state by the preceding theorem, while invertibility of
    the modular congruences and the supported-marginal correspondence give
    \begin{align}
        \dim_{\C}\range L
         & =\dim_{\C}\range\Phi_\rho
        =\operatorname{OSR}(\rho_{AB}).
        \notag
    \end{align}
    The rectangular Choi range-dimension estimate now proves the claim.
\end{proof}

\begin{theorem}[Real powers under unitary conjugation]
    \label{thm:rpow_conj_unitary}
    \lean{Matrix.rpow_conj_unitary}
    \leanok
    Let $A\succeq0$, let $s\in\R$, and let $U$ be unitary. Then
    \begin{align}
        (UAU^\dagger)^s & =UA^sU^\dagger.
        \notag
    \end{align}
\end{theorem}

\begin{proof}\leanok
    Apply Lemma~\ref{lem:cfc_conj_unitary} to the function $x\mapsto x^s$.
\end{proof}

\begin{theorem}[Real powers of positive tensor products]
    \label{thm:rpow_kronecker_possemidef}
    \lean{Matrix.PosSemidef.rpow_kronecker}
    \leanok
    Let $A\succeq0$ and $B\succeq0$. For every $s\in\R$,
    \begin{align}
        (A\otimes B)^s & =A^s\otimes B^s.
        \notag
    \end{align}
\end{theorem}

\begin{proof}\leanok
    Apply
    Theorem~\ref{thm:multiplicative_functional_calculus_tensor_product} to
    $x\mapsto x^s$, using $(xy)^s=x^sy^s$ for $x,y\geq0$.
\end{proof}

\begin{theorem}[Inverse quarter-power from iterated square roots]
    \label{thm:rpow_neg_quarter_inv_sqrt_sqrt}
    \lean{Matrix.PosDef.rpow_neg_quarter_eq_inv_sqrt_sqrt}
    \leanok
    If $A\succ0$, then
    \begin{align}
        A^{-1/4} & =\left(\sqrt{\sqrt A}\right)^{-1}.
        \notag
    \end{align}
\end{theorem}

\begin{proof}\leanok
    The continuous functional calculus gives
    $\sqrt{\sqrt A}=A^{1/4}$. Positive definiteness makes this matrix
    invertible, and inversion changes the exponent from $1/4$ to $-1/4$.
\end{proof}

\begin{theorem}[Unitary invariance of the order-two functional]
    \label{thm:sandwiched_renyi_two_unitary_invariance}
    \lean{TNLean.sandwichedRenyiTwoTrace_conj_unitary}
    \leanok
    \uses{def:sandwiched_renyi_two_trace}
    Let $\omega\succeq0$, let $\rho$ be a square matrix of the same size, and
    let $U$ be unitary. Then
    \begin{align}
        Q_2(U\rho U^\dagger,U\omega U^\dagger)
         & =Q_2(\rho,\omega).
        \notag
    \end{align}
\end{theorem}

\begin{proof}\leanok
    \uses{thm:rpow_conj_unitary}
    Theorem~\ref{thm:rpow_conj_unitary} gives
    $(U\omega U^\dagger)^{-1/4}=U\omega^{-1/4}U^\dagger$.
    Substitute this identity into both sandwiching factors and use
    $U^\dagger U=1$ and cyclicity of the trace.
\end{proof}

\begin{theorem}[Reindexing invariance of the order-two functional]
    \label{thm:sandwiched_renyi_two_submatrix_equiv}
    \lean{TNLean.sandwichedRenyiTwoTrace_submatrix_equiv}
    \leanok
    \uses{def:sandwiched_renyi_two_trace}
    Let $\rho$ be a square matrix and let $\omega\succeq0$ have the same finite
    index set. Let $e$ be a bijection onto another finite index set. Then
    \begin{align}
        Q_2(\rho_{e^{-1},e^{-1}},\omega_{e^{-1},e^{-1}})
         & =Q_2(\rho,\omega).
        \notag
    \end{align}
\end{theorem}

\begin{proof}\leanok
    \uses{def:sandwiched_renyi_two_trace}
    Reindexing commutes with the continuous functional calculus, matrix
    multiplication, and the trace. Apply these three identities to the two
    inverse quarter-powers and the two copies of the sandwiched matrix.
\end{proof}

\begin{theorem}[Faithful marginal whitening bound]
    \label{thm:sandwiched_renyi_two_osr_faithful}
    \lean{TNLean.sandwichedRenyiTwoTrace_product_marginals_le_operatorSchmidtRank_of_marginals_posDef}
    \leanok
    \uses{def:sandwiched_renyi_two_trace}
    Let $\rho_{AB}\succeq0$ have positive definite marginals $\rho_A$ and
    $\rho_B$. Then
    \begin{align}
        Q_2(\rho_{AB},\rho_A\otimes\rho_B)
         & \leq\operatorname{OSR}(\rho_{AB}).
        \label{eq:mpdo_faithful_q2_osr}
    \end{align}
    No trace normalization is required.
\end{theorem}

\begin{proof}\leanok
    \uses{lem:partial_trace_product_isometry, thm:operator_schmidt_rank_local_isometry, thm:full_support_eigenbasis_whitened_choi_bound, thm:rpow_kronecker_possemidef, thm:rpow_neg_quarter_inv_sqrt_sqrt, thm:sandwiched_renyi_two_unitary_invariance}
    Choose a unitary $U$ that diagonalizes the first marginal and conjugate
    $\rho_{AB}$ by $U^\dagger\otimes1$. The two marginals become
    $\diag(p_i)$ and $\rho_B$, while positivity and operator-Schmidt rank are
    unchanged. The transpose in the Choi whitening convention becomes
    invisible only in this eigenbasis, since $\diag(p_i)^T=\diag(p_i)$.

    Theorems~\ref{thm:rpow_kronecker_possemidef} and
    \ref{thm:rpow_neg_quarter_inv_sqrt_sqrt} identify the negative
    quarter-power of the product marginal with the two whitening factors in
    Theorem~\ref{thm:full_support_eigenbasis_whitened_choi_bound}. That theorem
    bounds the squared whitened trace by the ordinary operator-Schmidt rank.
    Finally, Theorem~\ref{thm:sandwiched_renyi_two_unitary_invariance} and the
    invariance of operator-Schmidt rank under local unitaries return
    to the original basis.
\end{proof}

\begin{theorem}[Marginal-support whitening bound]
    \label{thm:sandwiched_renyi_two_osr_support}
    \lean{TNLean.sandwichedRenyiTwoTrace_product_marginals_le_operatorSchmidtRank}
    \leanok
    \uses{def:sandwiched_renyi_two_trace}
    Let $\rho_{AB}\succeq0$ be any finite-dimensional bipartite operator. Then
    \begin{align}
        Q_2(\rho_{AB},\rho_A\otimes\rho_B)
         & \leq\operatorname{OSR}(\rho_{AB}).
        \label{eq:mpdo_support_q2_osr}
    \end{align}
    No trace normalization or positive-dimension hypothesis is required. The
    statement includes the cases in which either marginal support has
    dimension zero.
\end{theorem}

\begin{proof}\leanok
    \uses{thm:sandwiched_renyi_two_osr_faithful, thm:sandwiched_renyi_two_support_compression, thm:product_marginals_kernel, thm:operator_schmidt_rank_marginal_support_compression, thm:faithful_marginals_support_compression}
    Compress both factors simultaneously to the supports of $\rho_A$ and
    $\rho_B$. The compressed marginals are positive definite by
    Theorem~\ref{thm:faithful_marginals_support_compression}, even when a
    support coordinate space has dimension zero. The order-two functional is
    unchanged by
    Theorem~\ref{thm:sandwiched_renyi_two_support_compression}, and the
    ordinary operator-Schmidt rank is unchanged by
    Theorem~\ref{thm:operator_schmidt_rank_marginal_support_compression}.
    Apply Theorem~\ref{thm:sandwiched_renyi_two_osr_faithful} to the compressed
    operator.
\end{proof}

The sandwiched R\'enyi comparison also uses the following unnormalized
trace term.  For trace-one $\rho$ and $\alpha>1$, it is the quantity inside the
logarithm of the sandwiched R\'enyi divergence
\cite[Equation~(3)]{Beigi2013Sandwiched}
\cite[Definition~2]{MullerLennert2013Renyi} whenever
$\rho,\omega\geq0$ and $\ker\omega\subseteq\ker\rho$.  This includes singular
$\omega$, with negative powers interpreted as the generalized inverse on its
support.  In this regime, finiteness of the divergence requires the support
inclusion.  The present application uses $1<\alpha\leq2$; the order-one case is
the Umegaki endpoint.  The statements below concern only this algebraic
expression.  They are not assertions of
\cite[Proposition~4.5]{Cirac2016MPDO_arXiv}, and they do not establish
continuity or monotonicity in the order, interpolation between the endpoints,
differentiability at order one, a limit to Umegaki relative entropy, a
logarithmic divergence, or a comparison with relative entropy.

\begin{definition}[Total sandwiched R\'enyi trace]
    \label{def:sandwiched_renyi_trace}
    \lean{TNLean.sandwichedRenyiTrace}
    \leanok
    For square complex matrices $\rho$ and $\omega$ and every $\alpha\in\R$,
    set
    \begin{align}
        \widetilde Q_\alpha(\rho\Vert\omega)
         & =\re\tr\!\left[
            \left(\omega^{(1-\alpha)/(2\alpha)}
            \rho\,\omega^{(1-\alpha)/(2\alpha)}\right)^\alpha
            \right].
        \label{eq:sandwiched_renyi_trace}
    \end{align}
    Real powers are defined by continuous functional calculus, with negative
    powers equal to zero on the zero eigenspace.  Thus
    $\widetilde Q_\alpha$ is total even when $\alpha=0$ or $\omega$ is singular.
    For $\alpha>1$, on positive semidefinite inputs satisfying
    $\ker\omega\subseteq\ker\rho$, this convention gives the finite
    sandwiched R\'enyi trace term.  In this regime, if the support inclusion
    fails, the totalized value remains finite but is not the divergence trace
    term.
\end{definition}

\begin{theorem}[Nonnegativity on the faithful domain]
    \label{thm:sandwiched_renyi_trace_nonneg}
    \lean{TNLean.sandwichedRenyiTrace_nonneg}
    \leanok
    \uses{def:sandwiched_renyi_trace}
    If $\rho\geq0$ and $\omega>0$, then, for every $\alpha\in\R$,
    \begin{align}
        0 & \leq\widetilde Q_\alpha(\rho\Vert\omega).
        \notag
    \end{align}
\end{theorem}

\begin{proof}\leanok
    \uses{def:sandwiched_renyi_trace}
    Put $q=\omega^{(1-\alpha)/(2\alpha)}$.  Continuous functional calculus
    gives $q\geq0$, hence $q\rho q\geq0$ and $(q\rho q)^\alpha\geq0$.
    The trace of the last matrix is real and nonnegative.
\end{proof}

\begin{theorem}[Order-one sandwiched trace]
    \label{thm:sandwiched_renyi_trace_one}
    \lean{TNLean.sandwichedRenyiTrace_one}
    \leanok
    \uses{def:sandwiched_renyi_trace}
    If $\rho\geq0$ and $\omega>0$, then
    \begin{align}
        \widetilde Q_1(\rho\Vert\omega) & =\re\tr(\rho).
        \notag
    \end{align}
\end{theorem}

\begin{proof}\leanok
    \uses{def:sandwiched_renyi_trace}
    At $\alpha=1$ the two powers of $\omega$ have exponent zero, while the
    remaining real power of $\rho$ has exponent one.  Substitution into
    (\ref{eq:sandwiched_renyi_trace}) gives the identity.
\end{proof}

\begin{theorem}[Order-two sandwiched trace]
    \label{thm:sandwiched_renyi_trace_two}
    \lean{TNLean.sandwichedRenyiTrace_two}
    \leanok
    \uses{def:sandwiched_renyi_trace, def:sandwiched_renyi_two_trace}
    If $\rho\geq0$ and $\omega>0$, then
    \begin{align}
        \widetilde Q_2(\rho\Vert\omega) & =Q_2(\rho,\omega).
        \notag
    \end{align}
\end{theorem}

\begin{proof}\leanok
    \uses{def:sandwiched_renyi_trace, def:sandwiched_renyi_two_trace}
    At $\alpha=2$ the sandwiching exponent is $-1/4$.
    Since $\omega^{-1/4}\rho\,\omega^{-1/4}$ is positive semidefinite, its
    continuous-functional-calculus square is its ordinary matrix square.
\end{proof}

\begin{theorem}[Mutual information from operator-Schmidt rank]
    \label{thm:mutual_information_le_log_operator_schmidt_rank}
    \lean{Entropy.mutualInformation_le_log_operatorSchmidtRank}
    \leanok

    Let $\rho_{AB}\succeq0$ be a finite-dimensional bipartite operator with
    $\tr\rho_{AB}=1$. Then
    \begin{align}
        I(A:B)_\rho
        \leq \log\operatorname{OSR}(\rho_{AB}).
        \notag
    \end{align}
    No positive-dimension or faithful-marginal hypothesis is required. The
    operator-Schmidt rank is the ordinary operator-Schmidt rank.
\end{theorem}
\begin{proof}\leanok
    \uses{thm:product_marginals_kernel, thm:sandwiched_renyi_two_submatrix_equiv, thm:relative_entropy_q2_support, thm:sandwiched_renyi_two_nonneg, thm:sandwiched_renyi_two_osr_support, thm:operator_schmidt_rank_pos}
    Set $\omega=\rho_A\otimes\rho_B$. The marginals are positive semidefinite,
    $\tr\omega=1$, and Theorem~\ref{thm:product_marginals_kernel} gives
    $\ker\omega\subseteq\ker\rho_{AB}$. Reindex the product basis by a single
    finite coordinate set. Lemma~\ref{lem:relative_entropy_submatrix_equiv} and
    Theorem~\ref{thm:sandwiched_renyi_two_submatrix_equiv} preserve the two
    functionals, so Theorem~\ref{thm:relative_entropy_q2_support} applies.
    Theorem~\ref{thm:relative_entropy_product_marginals} identifies its left
    side with mutual information. Thus, with
    $q=Q_2(\rho_{AB},\omega)$,
    \begin{align}
        I(A:B)_\rho
          & =D(\rho_{AB}\Vert\omega)
        \leq\log q,
        \notag                                       \\
        0 & \leq q\leq\operatorname{OSR}(\rho_{AB}).
        \notag
    \end{align}
    The second inequality is
    Theorem~\ref{thm:sandwiched_renyi_two_osr_support}. Since
    $\tr\rho_{AB}=1$, the operator is nonzero, and
    Theorem~\ref{thm:operator_schmidt_rank_pos} gives
    $\operatorname{OSR}(\rho_{AB})>0$. This positive integer is at least one,
    so its logarithm is nonnegative. If $q=0$, the totalized convention
    $\log0=0$ gives $I(A:B)_\rho\leq0\leq\log\operatorname{OSR}(\rho_{AB})$.
    If $q>0$, monotonicity of the logarithm gives
    $\log q\leq\log\operatorname{OSR}(\rho_{AB})$.
\end{proof}
